% CV d'Issa Aziz KARAMBIRI
% Compiler avec pdflatex (Overleaf : pdfLaTeX)
% Version stage : laisser \stagetrue. Version simple (sans recherche de stage) : mettre \stagefalse.
\documentclass[10pt]{article}
\usepackage[utf8]{inputenc}
\usepackage[T1]{fontenc}
\input{glyphtounicode}\pdfgentounicode=1
\usepackage[french]{babel}
\usepackage[lining]{ebgaramond}
\linespread{1.0}
\AtBeginDocument{\fontsize{10.6}{12.9}\selectfont}
\usepackage[a4paper,top=1.0cm,bottom=0.9cm,left=1.6cm,right=1.6cm]{geometry}
\usepackage{xcolor}
\usepackage{titlesec}
\usepackage{enumitem}
\usepackage[hidelinks,pdfauthor={Issa Aziz KARAMBIRI},pdftitle={CV Issa Aziz KARAMBIRI - Élève-ingénieur data science, IA et finance},pdfsubject={CV},pdfkeywords={data science, intelligence artificielle, machine learning, data analyse, data engineering, finance, Python, SQL, stage},pdflang={fr}]{hyperref}
\usepackage{tabularx}

\definecolor{gris}{HTML}{6B7280}
\definecolor{trait}{HTML}{C9CED6}
\newif\ifstage
\stagetrue
\ifdefined\cvsimple\stagefalse\fi
\pagestyle{empty}
\setlength{\parindent}{0pt}
\frenchsetup{ItemLabels=\textbullet}

% Titres de section : gras, trait gris dessous
\titleformat{\section}{\large\bfseries}{}{0pt}{}[\vspace{-0.45ex}{\color{trait}\titlerule[0.6pt]}]
\titlespacing*{\section}{0pt}{1.1ex}{0.7ex}

% Puces petites et grises
\setlist[itemize]{leftmargin=1.3em,itemsep=0pt,topsep=1pt,parsep=0pt,label={\color{gris}\small\textbullet}}

% \entree{Titre}{détail en gris}{date}
\newcommand{\entree}[3]{\par\noindent\textbf{#1}\ifx&#2&\else\ {\color{gris}-- #2}\fi\ \hfill{\color{gris}#3}\par\vspace{1pt}}

\begin{document}

{\fontsize{24}{26}\selectfont\bfseries Issa Aziz KARAMBIRI}\par\vspace{3pt}
Élève-ingénieur en mathématiques appliquées et informatique, ENSIIE\par
\ifstage\textbf{Recherche un stage de 3 à 4 mois à partir de mai 2027 en data science, IA, data analyse ou finance}\par\fi\vspace{2pt}
Évry-Courcouronnes (91000) \quad\textperiodcentered\quad +33 7 61 20 07 82 \quad\textperiodcentered\quad \href{mailto:issaaziz.karambiri@ensiie.eu}{issaaziz.karambiri@ensiie.eu}\par
\href{https://issakarambiri.com}{Portfolio : issakarambiri.com} \quad\textperiodcentered\quad \href{https://github.com/KARAMBIRI65}{github.com/KARAMBIRI65} \quad\textperiodcentered\quad \href{https://www.linkedin.com/in/karambiri-ma}{linkedin.com/in/karambiri-ma}\par

\section*{Profil}
\ifstage
Élève-ingénieur en 2\textsuperscript{e} année à l'ENSIIE (Master 1), je cherche un stage où l'on me confie de vrais sujets, quitte à sortir de ma zone de confort. Mes projets m'ont surtout montré tout ce qu'il me reste à apprendre, et c'est exactement ce qui me motive. Je pose des questions, j'essaie, je me trompe parfois, et je progresse vite. Je veux rejoindre une équipe exigeante, apprendre à son contact et prendre des responsabilités au fil du stage.
\else
Élève-ingénieur en 2\textsuperscript{e} année à l'ENSIIE (Master 1), je me forme en mathématiques appliquées et en informatique avec une orientation data, IA et finance. Mes projets m'ont surtout montré tout ce qu'il me reste à apprendre, et c'est exactement ce qui me motive. Je pose des questions, j'essaie, je me trompe parfois, et je progresse vite. J'aime les sujets concrets, les équipes exigeantes et prendre des responsabilités.
\fi

\section*{Expériences}
\entree{Stagiaire data science, OptiTrade}{ENTSI, Oujda (Maroc)}{mars -- mai 2025}
\begin{itemize}
  \item Application web d'aide à l'investissement : 300\,000 lignes de données boursières au départ, 246\,437 après nettoyage.
  \item Quatre modèles comparés (régression linéaire, Random Forest, XGBoost, LSTM) avec MSE, RMSE et R², pour conseiller d'acheter, vendre ou garder une action.
  \item Python, Django REST, React, PostgreSQL, Scikit-learn, TensorFlow.
\end{itemize}
\vspace{3pt}
\entree{Stagiaire développeur web, GESCOM}{ENTSI, Oujda (Maroc)}{juin -- juil. 2024}
\begin{itemize}
  \item Application de gestion commerciale (clients, produits, ventes, achats, rôles utilisateurs) en PHP, MySQL, Bootstrap et jQuery.
  \item Optimisation des requêtes SQL les plus lentes.
\end{itemize}

\section*{Projets}
\entree{Maintenance prédictive des moteurs d'avion}{machine learning, données NASA}{\href{https://github.com/KARAMBIRI65/maintenance-predictive-moteurs-avion}{GitHub}}
\begin{itemize}
  \item Prévision du nombre de vols restants avant la panne (XGBoost, erreur typique de 11 vols) et seuil d'inspection choisi pour minimiser le coût total.
  \item 25 pannes imminentes détectées sur 25, avec seulement 5 fausses alertes sur 75 moteurs sains.
\end{itemize}
\vspace{3pt}
\entree{VIGIE, assistant anti-arnaques pour les seniors}{Challenge Projet, 4 jours}{}
\begin{itemize}
  \item Équipe de 5 étudiants de 3 écoles (ENSIIE, Télécom SudParis, IMT-BS). Mon rôle : responsable technique et prototypage.
  \item Détection des liens frauduleux, alerte avant le clic et écrans lisibles pour une personne de 74 ans.
\end{itemize}
\vspace{3pt}
\entree{Détection de fraude bancaire}{données très déséquilibrées}{\href{https://github.com/KARAMBIRI65/detection-fraude-bancaire}{GitHub}}
\begin{itemize}
  \item Fraudes représentant 0,17\,\% des transactions. XGBoost réglé en raisonnant en euros économisés (PR-AUC 0,88).
  \item Pertes liées à la fraude réduites de 77\,\%, avec 87\,\% d'alertes justes.
\end{itemize}
\vspace{3pt}
\entree{Détection de déchets sur les plages}{deep learning}{janv. -- mars 2025}
\begin{itemize}
  \item Images annotées à la main (6 classes), augmentation des données, entraînement de YOLOv8 (PyTorch) sur GPU et interface OpenCV.
  \item Précision 89,4\,\%, rappel 89,5\,\%, F1-score 89,5\,\%.
\end{itemize}
\vspace{3pt}
\textbf{Projets d'école :} modèle de Black-Scholes en C++ (prix d'options européennes), machine virtuelle OCaml en C (bytecode, pile, mémoire), mini-SGBD relationnel en OCaml.\par
{\color{gris}Tous mes projets : \href{https://issakarambiri.com}{issakarambiri.com}}

\section*{Formation}
\entree{École Nationale Supérieure d'Informatique pour l'Industrie et l'Entreprise (ENSIIE)}{}{depuis sept. 2025}
Cycle ingénieur, 2\textsuperscript{e} année (Master 1), Évry-Courcouronnes. Probabilités, processus stochastiques, optimisation, recherche opérationnelle, machine learning, deep learning, bases de données SQL et NoSQL.\par\vspace{3pt}
\entree{DUT Technologique, mention Très Bien}{École Supérieure de Technologie d'Oujda (Maroc)}{2023 -- 2025}

\section*{Compétences et intérêts}
\textbf{Langages :} Python, SQL, C++, C, R, OCaml, Java, PHP, HTML, CSS, JavaScript\par
\textbf{Data et IA :} Pandas, NumPy, Scikit-learn, XGBoost, TensorFlow, PyTorch, YOLOv8, OpenCV, SHAP\par
\textbf{Outils :} PostgreSQL, Django, Git, VS Code, Jupyter, Google Colab, LaTeX, Power BI, Excel\par
\textbf{Langues :} français (langue maternelle), anglais (B2)\par
\textbf{Centres d'intérêt :} lecture, musique, cuisine (toujours à la recherche de nouvelles recettes)

\section*{Engagement et distinctions}
\entree{Ambassadeur, YITU Association}{soutien scolaire et ateliers CV pour des étudiants}{depuis août 2026}
\entree{Chargé académique, AEBM}{coordination des activités académiques de l'association}{2024 -- 2025}
\entree{Concours d'art oratoire}{lauréat (CETASSO, Oujda, 2024), finaliste (AEBM, 2025)}{2024 -- 2025}
\entree{Prix et attestations d'excellence}{moyenne annuelle supérieure à 16/20 (CETASSO, AEBM, UEEMO)}{2023 -- 2024}

\end{document}
